\documentclass[12pt]{article}
\usepackage[utf8]{inputenc}
\usepackage[T1]{fontenc}
\usepackage{hyperref}
\title{Tunni Language Resources for K-12 Curriculum Development}
\author{Compiled by Claude Sonnet 5 (claude-sonnet-5) for Andr\'as Kornai}
\date{}
\begin{document}
\maketitle

\section{Introduction}
Tunni (Af-Tunni; ISO 639-3 tqq) is a Southern Somali variety, usually grouped with the Digil branch of Somali. It is spoken mainly in the Lower Shabelle, Middle Juba and Lower Juba regions of Somalia. Sources disagree on whether it is a separate language or a Somali dialect. Speaker counts also vary, from about 20,000 in older estimates to about 34,000 on Wikipedia. Some Tunni in Brava speak Bravanese (Chimini) rather than Af-Tunni. The search found one main grammar and some general reference pages, but no educational materials.

\section{Primary Educational Resources}
No primers, textbooks or school programmes were found in the search results.

\section{Linguistic and Reference Materials}
\begin{itemize}
\item Mauro Tosco, \emph{Af Tunni: Grammar, Texts and Glossary of a Southern Somali Dialect} (K\"oln: R\"udiger K\"oppe, 1997), the main description. It covers phonology, morphology, syntax, texts and a glossary: \url{https://www.koeppe.de/titel_print_af-tunni}
\item Tunni language article on Wikipedia: \url{https://en.wikipedia.org/wiki/Tunni_language}
\item Tunni people article on Wikipedia: \url{https://en.wikipedia.org/wiki/Tunni}
\item Endangered Languages Project entry: \url{https://endangeredlanguages.com/lang/5053}
\item A University of Rome Roma Tre repository record, which appears to be a thesis or related work; I did not confirm its contents: \url{https://arcadia.sba.uniroma3.it/handle/2307/2761?mode=full}
\end{itemize}

\section{Religious and Community Resources}
No Tunni scripture or community organisations were found in the search results.

\section{Digital and Audio Resources}
No Tunni audio or digital learning resources were found in the search results.

\section{Curriculum Development Context}
Documentation is limited to Tosco's grammar and general reference pages. Somali is the standard language of education in Somalia, and the relation between Tunni and Somali, and between Tunni and the Digil/Maay varieties, should be settled with the community before developing materials. The grammar and glossary could serve as a basis for an orthography and word lists.
\end{document}